% JETSPIN perturbations documentation

\section{Perturbations at the nozzle\label{sec:Perturbations-at-the}}

The spinneret nozzle is represented by a single mass-less point of
charge $\bar{q}$ fixed at $x=0$ (nozzle bead). Its charge $\bar{q}$
is assumed equal to the mean charge value of the jet beads. Such charged
point can be also interpreted as a small portion of jet which is fixed
to the nozzle. In JETSPIN it is possible to add small perturbations
to the $y_{n}$ and $z_{n}$ coordinates of the nozzle bead in order
to model fast mechanical oscillations of the spinneret\cite{coluzza2014ultrathin}.
Given the initial position of the nozzle

\begin{subequations}

\begin{equation}
y_{n}=A\cdot\cos\left(\varphi\right)
\end{equation}

\begin{equation}
z_{n}=A\cdot\sin\left(\varphi\right),
\end{equation}

\end{subequations}

the equations of motion for the nozzle bead are

\begin{subequations}

\begin{equation}
\frac{dy_{n}}{dt}=-\omega\cdot z_{n}
\end{equation}

\begin{equation}
\frac{dz_{n}}{dt}=\omega\cdot y_{n},
\end{equation}

\end{subequations}

where $A$ denotes the amplitude of the perturbation, while $\omega$
and $\varphi$ are its frequency and initial phase, respectively.
